\subsection{ODS và phân tích dữ liệu nguồn}
Dữ liệu từ hai hệ thống nguồn trước hết được chép nguyên trạng vào tầng staging (11 bảng), mỗi dòng gắn mã đợt nạp và thời điểm nạp để truy vết. Có hai chế độ trích xuất: toàn phần (xoá và chép lại) và tăng dần (chỉ lấy dòng mới dựa trên watermark: ngày mua của đơn, ngày liên hệ đầu của MQL, ngày chốt của deal; các bảng con lấy theo đơn mới). Từ staging, các thủ tục SQL làm sạch dữ liệu vào tầng ODS. Bảng toạ độ địa lý 1.000.163 điểm được gom về 19.010 mã bưu chính, lấy median toạ độ và mode tên thành phố sau khi loại 42 điểm nằm ngoài phạm vi Brazil. Tên thành phố được chuẩn hoá chữ thường, bỏ dấu (73.440 điểm toạ độ bị đổi). Cột sai chính tả \code{*\_\allowbreak{}lenght} được đổi tên; 610 sản phẩm không có danh mục được gán \code{unknown}; 13 sản phẩm thuộc hai danh mục thiếu bản dịch được bổ sung bằng bảng seed. Với đơn hàng, thay vì xoá, hệ thống gắn cờ lỗi: 129 trong 45.430 đơn có mốc thời gian sai thứ tự hoặc thiếu ngày (phổ biến nhất là ngày bàn giao vận chuyển trước ngày duyệt, 84 đơn). Các dòng không nạp được đều ghi vào bảng \code{etl\_\allowbreak{}log.rejected\_\allowbreak{}rows} kèm lý do; kiểm thử đối soát cho thấy với mọi bảng số dòng staging bằng số dòng ODS cộng số dòng bị loại.

\begin{table}[htbp]
\centering\small
\caption{Lỗi dữ liệu nguồn được phát hiện khi làm sạch và cách xử lý (số dòng ở trạng thái cuối)}
\label{tab:dq-errors}
\begin{tabular}{p{6.2cm}rp{5.5cm}}
\toprule
\textbf{Lỗi} & \textbf{Số dòng} & \textbf{Xử lý} \\
\midrule
Bàn giao vận chuyển trước khi duyệt thanh toán & 1.359 & Gắn cờ dq\_flags, giữ đơn; khoảng thời gian âm đặt NULL \\
Bàn giao vận chuyển trước ngày mua & 166 & Gắn cờ dq\_flags, giữ đơn; khoảng thời gian âm đặt NULL \\
Giao khách trước khi bàn giao vận chuyển & 23 & Gắn cờ dq\_flags, giữ đơn; khoảng thời gian âm đặt NULL \\
Đơn đã giao nhưng không có ngày duyệt & 14 & Gắn cờ dq\_flags, giữ đơn; khoảng thời gian âm đặt NULL \\
Trạng thái delivered nhưng không có ngày giao & 8 & Gắn cờ dq\_flags, giữ đơn; khoảng thời gian âm đặt NULL \\
Có ngày giao nhưng trạng thái chưa là delivered & 6 & Gắn cờ dq\_flags, giữ đơn; khoảng thời gian âm đặt NULL \\
Đã giao nhưng không có ngày bàn giao & 2 & Gắn cờ dq\_flags, giữ đơn; khoảng thời gian âm đặt NULL \\
Giá trị thanh toán = 0 & 9 & Gắn cờ, chuẩn hoá (số kỳ 0 → 1) \\
Hình thức thanh toán not\_defined & 3 & Gắn cờ, chuẩn hoá (số kỳ 0 → 1) \\
Số kỳ trả góp = 0 (đặt thành 1) & 2 & Gắn cờ, chuẩn hoá (số kỳ 0 → 1) \\
Doanh thu tháng khai báo = 0 & 797 & Gắn cờ, giữ nguyên giá trị \\
Chốt deal trước ngày liên hệ đầu & 1 & Gắn cờ, giữ nguyên giá trị \\
coordinates outside Brazil bounding box & 42 & Loại khỏi ODS, ghi etl\_log.rejected\_rows \\
\bottomrule
\end{tabular}
\\[2pt]{\footnotesize Cờ chất lượng gắn trên đơn, thanh toán và deal; mỗi dòng có thể có nhiều cờ.}
\end{table}

\todo{Nhóm bổ sung phần ``Giới thiệu về ODS'' theo giáo trình (khái niệm, vị trí giữa nguồn và kho, khác biệt với kho) và liên hệ với Chương 1.}
